% ─────────────────────────────────────────────────────────────────────────────
\section{La marge de s\'ecurit\'e suppl\'ementaire}
% ─────────────────────────────────────────────────────────────────────────────

Le code permet d'ajouter une \textbf{marge de s\'ecurit\'e suppl\'ementaire}
via le param\`etre \py{marge\_securite} (entre 0 et 1).
Elle est appliqu\'ee \emph{apr\`es} tous les calculs EC5,
en multipliant tous les taux par $(1 + \text{marge})$~:

\begin{equation}
\text{taux\_effectif} = \text{taux}_{EC5} \times (1 + \text{marge\_securite})
\end{equation}

\begin{lstlisting}[style=python, caption={\fichier{moteur.py} --- application de la marge de sécurité (EF-026)}]
def _appliquer_marge(résultats: list[dict], marge: float) -> list[dict]:
    """Applique taux_effectif = taux × (1 + marge) sur tous les taux."""
    if marge == 0.0:
        return résultats   # optimisation : rien à faire si marge = 0

    champs_taux = [k for k in résultats[0].keys() if "taux" in k]
    for r in résultats:
        for k in champs_taux:
            v = r.get(k)
            if isinstance(v, float):
                r[k] = v * (1.0 + marge)
    return résultats
\end{lstlisting}

Exemple~: avec \py{marge\_securite = 0.15} et un taux calcul\'e de 0{,}85~:
\[
\text{taux\_effectif} = 0{,}85 \times 1{,}15 = 0{,}978 \leq 1{,}0
\quad \Rightarrow \text{ admis, mais conservateur}
\]